%auto-ignore
%      this ensures the arxiv doesn't try to start TeXing here.
%!TEX root = super_lattice_models_draft.tex
%      prev line helps TeXShop do the right thing

\section{Conclusion}
\label{sec:conclusion}

We have developed a comprehensive approach to finding and exploiting topological defects in two-dimensional classical lattice models, including their quantum-spin-chain limits. This approach requires that the models be defined using a fusion category, with the transfer matrix written in terms of projection operators best defined using fusion diagrams. Such models include some of the best-known models of lattice statistical mechanics, including the Ising, Potts and self-dual eight-vertex models. The category setup allows our knowledge of lattice topological defects to be expanded significantly, making straightforward the derivation of properties such as their behavior under branching and fusing.  Linear identities between partition functions with different defect configurations follow easily.

These identities have a number of elegant applications.
They give a very useful implementation of the venerable Kramers-Wannier duality, allowing us in Part~I to find new results for the even more venerable Ising model. Here we derived the analogous results for the Potts, clock and eight-vertex models, and showed how to generalize duality itself to a far-wider class of models. 
We computed universal quantities directly on the lattice using the defects, including the eigenvalues of Dehn twists and ratios of $g$-factors.
If the limit of the lattice model is a conformal field theory, then these results provide strong and exact constraints on what that CFT may be. 
Another consequence is the existence of exact degeneracies in the ground states and low-lying spectrum of non-critical models. 
Whereas the continuum description of such models is by an integrable gapped quantum field theory, the lattice models we analyze typically are not integrable. Self-duality provides a direct way of understanding why these degeneracies occur.

We showed that the lattice partition functions can be viewed as boundary conditions to the Turaev-Viro-Barrett-Westbury partition function, and gave a natural interpretation of the topological defects as elements in the Drinfeld center $\mcz(\mcc)$.
This was further corroborated by the computation of the Dehn twist eigenvalues which coincided with the topological spins of simple objects in $\mcz(\mcc)$.


Many generalizations await. 
One we have already worked out in detail is the implementation of the partition-function identities on the torus. 
We further find that when the continuum limit is a conformal field theory that we can provide an explicit map between the lattice partition function on the torus and the conformal field theory partition function for an arbitrary topological defect configuration on the lattice.
If the continuum limit is a CFT, we show how an arbitrary topological defect configuration maps onto characters of the Virasoro and extended algebras.

Another is to define (non-local) operators by terminating a topological defect somewhere in the bulk of the system. 
These terminations are naturally labeled by simple objects in the Drinfeld center. 
Indeed, we can view the termination as a boundary with the topology of a circle, and so has a natural action from the tube category.
We can study the behavior of the operator under 2$\pi$ twists, so as to find the lattice version of the topological spin. 
As with the Dehn twist eigenvalues, the lattice values are independent of system size, and so if the continuum limit is a CFT, they must match the conformal spins of the corresponding operators. 
We thus find a host of generalizations of fermion and parafermion operators, beyond the conserved currents described in \cite{Fendley2020}. 
In particular, here it is possible to analyze generally terminations that change the $\upsilon$ used on an edge. 
Moreover, the braiding coming from the Drinfeld center allows the results of  \cite{Fendley2020} to be implemented solely using fusion category data, so that the braiding used there is not required.

Even though our construction gives a host of lattice topological defects and corresponding dualities, it is not exhaustive. 
A very interesting generalization is to construct topological defects separating Morita-equivalent fusion categories.
A pair of fusion categories are Morita equivalent if there exists an invertible bimodule between them. 
Consequently they have the same Drinfeld center and so one expects their lattice partition functions to have similar properties. 
Indeed, one can use the invertible bimodule to construct a lattice topological defect between any two Morita equivalent fusion categories, thus generalizing the dualities studied in this paper.
We note that Morita-equivalent fusion categories may look wildly different, for example $\text{Rep}(G)$ and $\text{Vec}_G$ are Morita equivalent but have different objects, fusion rules, etc.
Consequently, the dualities give precise relations between partition functions which have very different microscopic degrees of freedom.


Another generalization of our construction gives orbifold defects. 
Orbifold defects again relate partition functions with different microscopic degrees of freedom, 
but unlike the Morita-equivalence dualities, orbifold defects are non-invertible.
These defects allow one to extend an ancient result \cite{Fendley1989} relating partition functions of seemingly distinct lattice models via an analog of the orbifold construction of conformal field theory. 
We illustrate the idea by considering one of the old examples, where the ABF models with $k$ even are related to the $D_n$ series of models  \cite{Pasquier1986} with $n=2+\tfrac{k}{2}$.
The $D_n$ Dynkin diagram gives the adjacency graph of the latter, replacing the $\mathcal{A}_{k+1}$ one from \eqref{Akadjacency}. For example, the $D_4$ model has the incidence diagram of the three-state Potts model, associated with the $\mathbb{Z}_3$ Tambara-Yamagami fusion category.
Since the relation of Boltzmann weights of the $D_n$ model with the corresponding $\mathcal{A}_{2n-3}$ model is a linear one \cite{Fendley1989}, it is natural to expect a topological defect separating the region of the two models.
Indeed, with some work the explicit form of the orbifold defects can be extracted from \cite{Fendley1989}.
In the $D_3=\mathcal{A}_3$ Ising case, such a defect is precisely the KW duality we described above. 
However, for higher $n$, our construction requires generalization. 
As with their continuum counterparts \cite{Frohlich2009}, there exists a topological ``orbifold'' defect that implements the lattice orbifold. 
Similarly to the Morita-equivalence defects, these orbifold defects can be parametrized by Frobenius algebra objects in $\mcc$.


\paragraph{Pre-acknowledgments} This work was done almost entirely in 2015. Part I appeared in January 2016, and we gave multiple talks describing the main results, e.g.\ \cite{Aasen2017,Fendley2018}. The only more recent results were a few small applications, mainly the Dehn-twist calculation in the clock models.  Over the five years we (fairly inexplicably) took to write up these results, a number of nice papers have appeared with some overlap. A partial list includes
\cite{Buican2017,freed2018,Vanhove2018,Belletete2018,thorngren2019,Wenjie2019,Chang_2019,lootens2019,Belletete2020}.

\paragraph{Acknowledgments} It is a pleasure to thank Jason Alicea, Kevin Walker, and Dominic Williamson for stimulating discussions.
DA was supported by a postdoctoral fellowship from the Gordon
and Betty Moore Foundation, under the EPiQS initiative, Grant GBMF4304.
PF was supported by EPSRC grants EP/S020527/1 and EP/N01930X (PF), while 
RM was supported by NSF DMR-1848336.

